% Generated draft: requires mathematical and source review.
\documentclass[11pt]{article}
\usepackage[margin=25mm]{geometry}
\usepackage{fontspec}
\setmainfont{Latin Modern Roman}
\usepackage{amsmath,amssymb,mathtools,unicode-math}
\setmathfont{Latin Modern Math}
\usepackage{enumitem,needspace,xurl,hyperref}
\setlength{\parindent}{0pt}
\setlength{\parskip}{5pt}
\setlength{\emergencystretch}{3em}
\newcommand{\N}{\mathbb{N}}
\newcommand{\Z}{\mathbb{Z}}
\newcommand{\Q}{\mathbb{Q}}
\newcommand{\R}{\mathbb{R}}
\newcommand{\C}{\mathbb{C}}
\newcommand{\F}{\mathbb{F}}
\newcommand{\A}{\mathbb{A}}
\newcommand{\PP}{\mathbb{P}}
\newcommand{\E}{\mathbb{E}}
\newcommand{\eps}{\varepsilon}
\newcommand{\dd}{\,\mathrm{d}}
\begin{document}
\section*{UnsolvedMath ID: 30000297}
\textbf{Source problem code:} catalogue entry
\subsection{Definitions and mathematical statement}
Using the cleaned statement's base field, let \(G\) be a noncommutative algebraic group over \(\C\). A compactification in the usual algebraic sense is a proper complex algebraic variety \(\overline G\) together with an open immersion
\[
j:G\hookrightarrow\overline G
\]
whose image is dense. The source requests constructions of such compactifications for nonabelian groups, without specifying connectedness, linearity, reductivity, or a particular group.

The extracted question does not impose smoothness, projectivity, a normal-crossings boundary, a prescribed orbit stratification, or equivariance. For example, a two-sided equivariant compactification would additionally extend the action \((a,b)\cdot g=agb^{-1}\) of \(G\times G\), but that is a stronger, separately specified target. Nor does compactification by itself require the multiplication and inversion maps to extend as a group structure on \(\overline G\). These different requirements cannot be conflated.

The supplied assessment checked 2026-08-21 reports wonderful compactifications in semisimple adjoint and broad symmetric settings and calls the unrestricted extracted request broader than those constructions. It does not reproduce the properties intended of the compactifications in the open question. Therefore that partial-status assessment is not sufficient to identify a unique unresolved theorem about bare existence as defined above; the requested extra geometry, if any, must be recovered.

The raw original wording ends ``over C times.'' The supplied correction explicitly chooses the complex base field \(\C\), which is the formulation retained here. The raw phrase could carry lost notation or context involving \(\C^\times\); its correction is supplied curation rather than an independently verified reconstruction.
\subsection{Short English statement}
Construct compactifications of noncommutative complex algebraic groups with the properties intended in the source. The extracted wording omits the geometric or equivariant requirements that would distinguish the requested programme from bare compactification.
\subsection{Sources}
\begin{itemize}
\item Local UnsolvedMath registry, permanent ID 30000297, catalogue entry, \emph{Compactifications of Nonabelian Algebraic Groups}.
\item \emph{Analysis and Geometric Singularities} (2006), Oberwolfach Reports, pp. 2187--2240, DOI \url{https://doi.org/10.4171/owr/2005/39}.
\item Supplied reference labeled ``wonderful compactifications'': \url{https://arxiv.org/abs/1003.1704}; no full author list or exact theorem is reproduced.
\item Supplied notation-correction claim and literature triage checked 2026-08-21; the latter reports the semisimple-adjoint and symmetric partial settings.
\end{itemize}
\end{document}
